% chapters/03/section-01.tex
\section{Dari Teori ke Praktik}
\label{sec:teori-ke-praktik}

Bab~II membekali Anda dengan kerangka berpikir tentang waktu, ruang, fisika,
dan pengendalian objek. Bab ini memindahkan kerangka itu ke atas meja kerja:
menyiapkan perkakas, menyusun tempat bekerja, dan menulis program permainan
pertama yang benar-benar berjalan.

\subsection{Tiga Hasil yang Harus Anda Peroleh}

Bab ini selesai apabila tiga hal berikut sudah ada, bukan sekadar dipahami.

\begin{enumerate}
  \item \textbf{Lingkungan kerja yang siap.} Python beserta Pygame terpasang
        dan dapat dipanggil, dan Anda mampu memeriksa kesiapannya sendiri
        tanpa bantuan orang lain.
  \item \textbf{Program pertama yang berjalan.} Sebuah program Pygame yang
        membuka jendela, menampilkan teks informasi versi, dan menunggu sampai
        jendela ditutup.
  \item \textbf{Tempat kerja yang tertata.} Folder proyek dengan struktur yang
        jelas dan sebuah repositori Git yang menyimpan riwayat pekerjaan Anda
        dari minggu pertama.
\end{enumerate}

Ketiga hasil tersebut akan dipakai terus sampai akhir semester. Praktikum pada
minggu-minggu berikutnya menumpuk berkas baru pada struktur folder yang sama,
dan seluruh nilai portofolio diambil dari repositori yang Anda buat sekarang.
Karena itu, mengerjakan bagian ini dengan rapi jauh lebih menguntungkan
daripada mengerjakannya sekadar agar lewat.

\subsection{Konvensi Penulisan Kode pada Buku Ini}

Sebelum menulis kode pertama, ada satu konvensi yang perlu Anda ketahui agar
tidak membingungkan.

\begin{itemize}
  \item \textbf{Nama variabel, fungsi, dan kelas ditulis dalam bahasa Inggris},
        misalnya \code{screen}, \code{check_pygame}, atau \code{parse_arguments}.
        Pilihan ini diambil agar kode pada buku ini sejalan dengan dokumentasi
        resmi Pygame \mbox{[R3]} dan buku rujukan \mbox{[R1]}, \mbox{[R2]} yang
        seluruhnya berbahasa Inggris. Anda akan sering berpindah antara kode
        dan dokumentasi, dan kesamaan penamaan mempercepat pekerjaan itu.
  \item \textbf{Komentar, keterangan fungsi, dan keluaran program ditulis dalam
        bahasa Indonesia}, agar maksud setiap bagian kode dapat dipahami tanpa
        penerjemahan.
  \item \textbf{Istilah konsep pada teks buku ditulis dalam bahasa Indonesia},
        dan istilah asingnya dicetak miring pada kemunculan pertama. Jadi teks
        menyebut \istilah{selang waktu}, sedangkan kode memakai nama
        \code{delta_time} atau menampilkan kata yang sama pada keluaran
        program.
\end{itemize}

Konvensi penamaan ini berlaku pada seluruh buku, termasuk pada skrip C\# untuk
Unity nanti. Menurut pengalaman, kebiasaan ini juga berguna di luar
perkuliahan: kode yang dibaca orang lain---termasuk dosen, tim, atau
pewawancara kerja---lazim memakai penamaan berbahasa Inggris.

\subsection{Batas Pembahasan Bab Ini}

Bab ini tidak membahas tiga hal berikut, dan mengetahuinya akan mencegah Anda
mencari-cari bagian yang memang belum ada.

\begin{itemize}
  \item \textbf{Siklus permainan.} Program pada bab ini menggambar sekali lalu
        menunggu peristiwa. Perulangan yang memperbarui keadaan dan menggambar
        ulang secara terus-menerus baru dibahas pada Bab~IV.
  \item \textbf{Gerak dan fisika.} Objek yang bergerak, gravitasi, dan tabrakan
        mulai dibahas pada Bab~V sampai Bab~VII.
  \item \textbf{Pemasangan Unity.} Perkakas Blok~III dipasang tersendiri pada
        Bab~XII, karena kebutuhan perangkatnya berbeda dari Python.
\end{itemize}
